% Generated draft: requires mathematical and source review.
\documentclass[11pt]{article}
\usepackage[margin=25mm]{geometry}
\usepackage{fontspec}
\setmainfont{Latin Modern Roman}
\usepackage{amsmath,amssymb,mathtools,unicode-math}
\setmathfont{Latin Modern Math}
\usepackage{enumitem,needspace,xurl,hyperref}
\setlength{\parindent}{0pt}
\setlength{\parskip}{5pt}
\setlength{\emergencystretch}{3em}
\newcommand{\N}{\mathbb{N}}
\newcommand{\Z}{\mathbb{Z}}
\newcommand{\Q}{\mathbb{Q}}
\newcommand{\R}{\mathbb{R}}
\newcommand{\C}{\mathbb{C}}
\newcommand{\F}{\mathbb{F}}
\newcommand{\A}{\mathbb{A}}
\newcommand{\PP}{\mathbb{P}}
\newcommand{\E}{\mathbb{E}}
\newcommand{\eps}{\varepsilon}
\newcommand{\dd}{\,\mathrm{d}}
\begin{document}
\section*{UnsolvedMath ID: 30002412}
\textbf{Source problem code:} catalogue entry
\subsection{Definitions and mathematical statement}
An \emph{involution constraint} for an evolution system is a differential condition on its state that is propagated by the evolution: a sufficiently regular solution satisfying the condition initially continues to satisfy it for as long as the propagation property applies. In the usual interpretation for ideal magnetohydrodynamics (MHD), the magnetic field \(B\) obeys the solenoidal constraint
\[
\nabla\cdot B=0.
\]
This identifies the standard MHD constraint relevant to the question; the supplied statement does not explicitly specify the involution operator or exclude additional constraints. A Riemann solver computes an exact or approximate solution of an initial discontinuity problem for the governing equations, typically as part of a numerical method.

\textbf{Source question.} What is the role of the involution structure in numerical calculations, and can it influence the stability of Riemann solvers for ideal magnetohydrodynamics?

The requested investigation is to characterize the relationship between the treatment of involution constraints in numerical calculations and the stability of Riemann solvers for ideal MHD. In the standard solenoidal interpretation, this includes understanding how preservation or violation of \(\nabla\cdot B=0\) affects the relevant stability property. The source does not specify a spatial dimension, domain, boundary conditions, admissible states, equation of state, exact or approximate solver, numerical discretization, discrete divergence operator, perturbation class, or definition of stability. Consequently, it does not determine a single quantified mathematical proposition. These choices must be specified before a precise stability assertion can be formulated; none is imposed here.
\subsection{Short English statement}
Determine how involution constraints, conventionally including the divergence-free magnetic-field condition, affect numerical calculations and the stability of ideal-MHD Riemann solvers. The supplied question leaves the solver class and meaning of stability unspecified, so it remains a research question requiring further formulation.
\subsection{Sources}
\begin{itemize}
\item Local UnsolvedMath registry: supplied source record for permanent ID \(30002412\), source problem code \texttt{catalogue entry}, in \texttt{data/unsolvedmath/problems.json}. Title: \emph{Involution Structure and Stability of MHD Solvers}. The record supplies proposed year 2013 and no named proposer.
\item Supplied report citation: \emph{High-Resolution Mathematical and Numerical Analysis of Involution-Constrained PDEs} (2014), Oberwolfach Reports, pp.\ 2691--2747, DOI \texttt{10.4171/owr/2013/47}. \url{https://doi.org/10.4171/owr/2013/47}. The collection describes these reports as documenting mathematical workshops and their open-problem discussions.
\item Supplied literature source, labeled ``Oberwolfach Report'': \url{https://ems.press/content/serial-article-files/46478}.
\item Supplied literature source, labeled ``Kemm'': \url{https://www.math.ntnu.no/conservation/2010/018.html}.
\item Supplied literature source, labeled ``Gardiner--Stone'': \url{https://doi.org/10.1016/j.jcp.2004.11.016}.
\item Literature assessment dated exactly \texttt{2026-08-21}: the supplied curation labels the problem partially solved and reports that effects of involution violations on MHD stability have been documented and constraint-preserving solvers developed, while a complete general characterization of the requested role remains open. This is supplied dated curation derived from the OWR record, not an independently verified literature assessment. Its time-sensitive status and cited claims require independent verification.
\item The registry marks the statement as exact and records a manual review describing the extraction as self-contained, mathematically coherent, and free of LaTeX or report-layout artifacts. These are supplied provenance claims, not independent verification of the report or of mathematical completeness.
\end{itemize}
\end{document}
